\subsection{Nonabelian refinement tests at two and three}

\begin{statement}{D.27}{Proved}{Class-two obstructions, an alternate CM witness, and relative pro-prime lifts}
Keep $B,Q,N$ from D.25, write $\gamma_3N=[N,[N,N]]$, and put
\[
 K_3=\gamma_3N\,N^3,\qquad
 K_2=\gamma_3N\,[N,N]^2N^4,\qquad H_p=B/K_p.             \tag{D.27.1}
\]
These are characteristic refinements of the kernel $N$,
with
\[
 |H_3|=576\cdot3^{4753},\qquad
 |H_2|=576\cdot2^{4850}.                                \tag{D.27.2}
\]
The central layers over $N^{\rm ab}/p$ have dimensions
$4656$ at three and $4753$ at two.

At three, exact averaging over the normal subgroup
$O_2(Q)\simeq C_2^2\times Q_8$ reduces the central correction
system to a $135\times270$ matrix of rank 127. Among the
twenty-four \emph{chosen certificates} of D.25, four admit
class-two lifts with the recorded conjugators: two
CM-compatible pairs and the two extra pairs
\[
 (A,\beta)=\left(I,\begin{pmatrix}9&8\\4&9\end{pmatrix}\right),
 \quad
 \left(I,\begin{pmatrix}9&4\\8&9\end{pmatrix}\right)
 \pmod{12}.                                             \tag{D.27.3}
\]
Both extra lifts have zero central correction. Their fibre
products with $G_2$ give nonarithmetical finite shadows of
order $576\cdot2^{97}\cdot3^{4753}$.

At two the exact central correction matrix has rank 4737.
Allowing the multiplier $\lambda=1+72\eta$, $\eta\in\mathbb F_2$,
raises the rank to 4738. The fixed-conjugator test has two
survivors: the identity pair and the extra pair
\[
                  A=5I,\qquad\beta=I\pmod{12}.          \tag{D.27.4}
\]
Both use $\lambda=1$. Its fibre product with $G_3$ gives a
nonarithmetical finite shadow of order
$576\cdot2^{4850}\cdot3^{97}$.

There are actual infinite nonabelian extensions of these
three witnesses. Let $\mathcal P_p$ be the relative pro-$p$
completion of $B$ over $Q$, defined by retaining the full
pro-$p$ completion $N^{(p)}$ of $N$ as kernel. There is an
exact sequence
\[
 1\longrightarrow N^{(p)}\longrightarrow\mathcal P_p
 \longrightarrow Q\longrightarrow1,
 \qquad N^{(p)}\text{ free pro-}p\text{ of rank97}.        \tag{D.27.5}
\]
The two pairs (D.27.3) extend to automorphisms of
$G_2\times_Q\mathcal P_3$, and (D.27.4) extends to an
automorphism of $\mathcal P_2\times_QG_3$.
After one cusp gauge, each gives compatible finite GT
shadows on every finite characteristic refinement retaining
the indicated $G_6$ and CM data.
These are automorphisms of the specified profinite quotients.
The complete lower-conjugator freedom preserving each selected
$G_6$ action has dimension 52 at two and 34 at three.
Adjoining its correction columns raises the ranks to 4748
and 130, respectively, and leaves exactly the same survivors.
Their intersection consists only of the identity action.
Consequently \emph{all twenty-three nonidentity selected
$G_6$ actions in D.25 have no genuine full profinite lift}.
This includes all twelve selected extra-CM shadows, despite
their compatible lifts through D.26's entire homology tower.

There is, however, a different homology correction for the
first pair in (D.27.3) that lifts to both $H_2$ and $H_3$.
Its five exact mod-two obstruction forms vanish, the induced
kernel map has rank 97, and the corrected cusp word of length
470844 evaluates to the identity in the full Magnus model.
It has multiplier $\lambda=1$ and gives an automorphism of
\[
 H_2\times_Q H_3,
 \qquad |H_2\times_Q H_3|=576\cdot2^{4850}\cdot3^{4753}.
                                                               \tag{D.27.9}
\]
The same witness extends to $H_2\times_Q\mathcal P_3$,
with the identity action on $\mathcal P_3$ and prescribed
elliptic conjugators. A single cusp gauge gives compatible
nonarithmetical shadows on every finite characteristic
refinement retaining its $G_6$ and CM data. D.30 proves
that this particular recorded central representative has
no genuine full lift, jointly retaining its $G_6$ data,
and constructs a further explicit central replacement.
Genuine realization of the CM pair remains undecided.
\end{statement}
\begin{proof}
The characteristic quotients in (D.27.1) have explicit
free-group models. At three, use the class-two group
\[
 (v,z)(v',z')=(v+v',z+z'+\tfrac12v\wedge v'),
 \quad v\in\mathbb F_3^{97},\ z\in\bigwedge^2\mathbb F_3^{97}.
                                                               \tag{D.27.6}
\]
It has exponent three and is freely generated, in the
class-two exponent-three variety, by $(e_i,0)$.
Collection of commutators gives the asserted kernel
$\gamma_3N\,N^3$, dimension $97+\binom{97}{2}=4753$.
At two use the truncated noncommutative Magnus expansion
$a_i\mapsto1+X_i$ over $\mathbb F_2$, modulo degree three.
The multiplication is
\[
 (v,w)(v',w')=(v+v',w+w'+v(v')^{\mathsf T}).             \tag{D.27.7}
\]
The kernel over $v\in\mathbb F_2^{97}$ is the space of
symmetric $97\times97$ matrices: the diagonal entries are
$a_i^2$, and the off-diagonal basis entries are
$[a_i,a_j]$. They are independent central elements of order
two, proving the kernel description in (D.27.1) and the
order $2^{97+97+\binom{97}{2}}=2^{4850}$.
These normal forms also prove that the models are faithful.

Here is an exact evaluation of the required words in these
models. In the graph of D.25 retain its spanning tree and
orient every edge from its $s$-vertex to its $r$-vertex.
During evaluation of a word in $B$, keep its current image
$q\in Q$. Multiplication by $s$ changes the state to $qs$
and contributes no edge. Multiplication by $r$ appends
$e_q-e_{qr}$; multiplication by $r^{-1}$ appends
$e_q-e_{qr^2}$. Delete tree edges and read the remaining
oriented edges as free generators. For a closed $N$-word
this is exactly the tree-collapse isomorphism
$N\simeq\pi_1(N\backslash\mathcal T)\simeq F_{97}$.
The program constructs the $97$ basis words in $s,r$ from
the fundamental cycles and verifies that they evaluate
to the standard generators. Thus evaluating a cusp defect
uses actual group words, including all degree-two terms.

Let $M(q)$ be the exact matrix of $q$ on the cycle group
modulo $p$. The central action is
$z\mapsto M(q)zM(q)^{\mathsf T}$, on alternating matrices
at three and symmetric matrices at two. Left-multiplying
the two elliptic conjugators by central elements changes
the cusp defect by
\[
         (1-s)z_\beta+s(1-r)z_A.                        \tag{D.27.8}
\]
At three the uncorrected defect is computed from the word
$\beta s\beta^{-1}ArA^{-1}(sr)^{-1}$ in (D.27.6).
Its degree-one coordinate is zero by D.25. At two the
same word is evaluated in (D.27.7), with a symmetric
degree-two coordinate.

The subgroup $O_2(Q)$ has order 32. Indeed the unique lift
of $H=q(F)$ whose modulo-four component lies in $V_4$
identifies $H$ with $V_4\times\mathrm{SL}_2(\mathbb F_3)$;
$V_4\times Q_8$ is normal in $Q$. Averaging over it at
three is an exact idempotent. Its image on the central
layer has dimension 135. This can be computed both from
the explicit projector and from the integral graph
character
\[
 \chi_C(q)=576\,[q=1]-\#(Q/\langle s\rangle)^q
                -\#(Q/\langle r\rangle)^q+1,
 \quad
 \chi_{\wedge^2C}(q)=\tfrac12(\chi_C(q)^2-\chi_C(q^2)).
\]
The averaged integral lattice is a direct summand over
$\mathbb Z_{(3)}$, so its rank remains135 modulo three.
The program selects an explicit basis of that image and
computes the rank 127 matrix in (D.27.8).

Averaging loses no obstruction. If $P$ is the averaging
operator, then $z-Pz=\frac1{32}\sum_{g\in O_2(Q)}(1-g)z$.
Writing each $g$ as a word in $s,r$ telescopes each
$(1-g)z$ into $(1-s)a+(1-r)b$. The identity
\[
 (1-r)b=s(1-r)b+(1-s)(1-r)b
\]
converts this to (D.27.8). The program performs that
telescoping construction, reconstructs the central
correction in all 4656 coordinates, and checks the full
equation. The resulting four successful certificates
are retained in \texttt{d\_cm\_class2.json}.

At two the program uses exact integer bit vectors for all
4753 coordinates. The multiplier must retain
$m=0\pmod{36}$ at $G_6$. The two possibilities in the new
quotient are $\lambda=1+72\eta$ modulo 144. Since $t^{12}\in N$
and $N/K_2$ has exponent four,
$t^{72}=(t^{12})^2$ in $H_2$; its central coordinate is
$c c^{\mathsf T}$, where $c=[t^{12}]\bmod2$.
Appending this single column gives the rank 4738 augmented
system. For every successful solution the code checks the
complete4753-coordinate equation. It stores the exact
central coefficients and multiplier bit in
\texttt{d\_cm\_class2\_two.json}.

We now account for \emph{all} conjugator choices retaining
a selected $G_6$ action. Equation (D.25.6) forces their
images $q(A),q(\beta)\in Q$: the difference of two
conjugators inducing the same finite action centralizes
the corresponding elliptic generator, and the indicated
centralizers have trivial image in $Q$. The remaining
differences therefore lie in the kernel $N$.
At two, right-multiplication $\beta\mapsto\beta c$ can
preserve the $G_2$ action precisely when the degree-one
class of $c$ lies in $(N^{\rm ab}/2)^s$, a space of
dimension 52. For each basis class the program evaluates
the actual word $[c,s]$ in the Magnus model and adds its
central defect, transported by $q(\beta)$, to (D.27.8).
Changes of lift of $c$ are already included among its
central columns. The analogous right freedom of $A$
contributes no further obstruction: every order-three
centralizer element lifts through the two-kernel by
the coprime conjugacy argument proved below. At three
the roles are reversed. The space $(N^{\rm ab}/3)^r$
has dimension 34; its actual defects $[c,r]$, transported
by $s q(A)$, are appended to the averaged system.
The order-two centralizer lifts through the three-kernel.
The resulting exact ranks are 4748 and 130. These
calculations exhaust the lower-conjugator freedom,
and their saved certificates show the survivor sets
asserted in the statement.

Here is why their disjointness away from the identity
proves the stated genuine nonliftability. A genuine GT$_0$
action on the completed braid quotient induces an
automorphism of
$\widehat B=C_2\amalg C_3$ with its specified cusp power.
The torsion theorem of Herfort--Ribes,
\cite[Theorem A, p.~155; Theorem A$'$, pp.~157--159]{A_HerfortRibes1985},
makes the images of $s,r$ conjugates of $s,r$ respectively:
their orders distinguish the two factors, and the given
action on $B^{\rm ab}=C_6$ distinguishes $r$ from $r^{-1}$.
The induced identity on $S_3$ lets the conjugators be
chosen in $\widehat F$: multiply a conjugator by an
appropriate power of its elliptic generator to remove
its image in the corresponding $S_3$ centralizer.
Undo the known D.25 cusp gauge by the inner automorphism
$\operatorname{Int}(x^{-k})$. Since $x=u^2$ centralizes
the cusp $u$, the prescribed cusp power is retained.
The resulting automorphism induces the selected action
on $G_6$ and the identity on $Q$, hence preserves
$\widehat N$ and its characteristic subgroups defining
$H_2,H_3$. It must therefore give elliptic-conjugacy
lifts to \emph{both} of these quotients, with exactly
the conjugator freedom just exhausted. The multiplier
condition $m=0\pmod{36}$ is included: at two it is
$\lambda=1$ or73 modulo 144, while at three its change
acts trivially on the exponent-three kernel. Only the
identity selected action survives both systems.
This argument proves nonexistence of a genuine full
lift of each of the other twenty-three selected actions.
It does not rely on finite-quotient torsion conjugacy,
which need not follow from the full profinite theorem.

Every successful correction is an automorphism, with a
separate generation proof. The maps already induce the
identity on $Q$, preserve $N$, and retain the invertible
map on $N^{\rm ab}/p$ certified in D.25. The Frattini
subgroup of each finite $p$-kernel is its subgroup of
powers and commutators, so that map is onto the whole
$p$-kernel. Hence it is onto the extension and therefore
an automorphism. The two prime kernels have coprime orders;
the map from $B$ to their fibre product over $Q$ is onto
by the same property on $N$. This proves the stated
combined finite quotient orders and automorphisms.
The cusp gauge of D.25 supplies a commutator representative
and both hexagons, so the finite actions give actual shadows.

The alternate certificate is found and verified separately
in \texttt{d\_cm\_quadratic\_search.py}. The mod-two
degree-one cusp equation has a98-dimensional affine
solution space in 194 conjugator coordinates. The
orthogonal complement of the rank 4748 correction image
has dimension five. Applying those five linear
functionals to the actual degree-two word gives five
quadratic polynomials on the affine solution space.
The program computes them directly from the graph paths,
including the inverse-generator square terms. In the
fixed spanning-tree basis its successful homology masks are
\[
 b=\texttt{8ca9620ec63b26ee68fbb210}_{16},\qquad
 a=\texttt{60da0c00a2c8c0f44be10df0}_{16}.                \tag{D.27.10}
\]
They occur on the 48th tested affine vector; all five
obstructions vanish and the independently computed
kernel rank is 97. The exact central corrections and
right $s$-centralizer coefficients are saved in
\texttt{d\_cm\_quadratic\_search.json}. Substituting
them as explicit products of the 97 basis words, their
squares and commutators gives a470844-letter cusp word.
Direct evaluation returns $q=1$, $v=0$, and zero in
every one of the 4753 central coordinates. Thus existence
does not depend on random search success rates or solely
on the five projected equations. Retain the original
mod-three certificate, whose central correction is zero.
The coprime kernel product and the Frattini generation
argument already proved give (D.27.9).

We next prove the infinite relative lifts, rather than
inferring them from these finite calculations. Define
$\mathcal P_p$ as the inverse limit of $B/L$ over the
characteristic finite $p$-power-index subgroups $L$ of $N$.
Such subgroups are cofinal among the finite $p$-quotients
of the finitely generated free group $N$: take the
intersection of all normal subgroups of $p$-power index
bounded by a given integer.
This proves (D.27.5) and gives the canonical maps
$\widehat B\twoheadrightarrow\mathcal P_p\twoheadrightarrow G_p$.

For the two pairs (D.27.3), the recorded mod-three
conjugator $A$ is exactly the identity. The cusp equation
then says $\beta s\beta^{-1}=s$ in $G_3$, so the recorded
automorphism of $G_3$ is the identity. Its specific
centralizing $\beta$ lifts to an element of
$C_{\mathcal P_3}(s)$. To prove this, lift it arbitrarily;
its conjugate of $s$ and $s$ are two complements of order
two in an extension with pro-three kernel. They are
conjugate by that kernel. For a finite $p$-kernel and a
cyclic group of order prime to $p$, this conjugacy follows
by induction on the kernel: pass to an invariant central
elementary abelian subgroup and average the additive
cocycle there. Compactness gives the result for a
pro-$p$ kernel. Thus the identity automorphism of
$\mathcal P_3$ has the required elliptic conjugators and
reduces to the specified mod-three action.
Its fibre product with D.25's mod-two automorphism proves
the first assertion following (D.27.5).
Its fibre product with the alternate mod-two class-two
automorphism instead proves the asserted action on
$H_2\times_Q\mathcal P_3$.

For (D.27.4), the recorded mod-two $\beta$ is exactly the
identity; the cusp equation gives $ArA^{-1}=r$ in $G_2$.
The same coprime conjugacy proof, now with an order-three
complement and a pro-two kernel, lifts its prescribed $A$
to $C_{\mathcal P_2}(r)$. Combining the identity action on
$\mathcal P_2$ with the recorded mod-three action proves
the other assertion.

Finally choose a lift of each combined $\beta$ to
$\widehat F$, write its abelianization as $ax+by$, and
conjugate the quotient automorphism by $x^{b-a}$.
The resulting $f=s\beta's(\beta')^{-1}$ is the image of
an element of $\widehat F'$ and satisfies both hexagons
on every finite quotient. This single choice supplies
compatible shadows; a discrete representative agreeing
in the finite quotient and in $\widehat F^{\rm ab}/36$
gives the corresponding D.25 gauge. Their CM coordinates
are unchanged. Their nonarithmeticality therefore follows
from the same common-field restriction in D.25.

The strongest new inverse system retains the entire
free pro-three kernel and the explicit class-two
two-kernel. Its quotient maps are
$\widehat B\twoheadrightarrow H_2\times_Q\mathcal P_3
\twoheadrightarrow H_2\times_QH_3$.
Higher two-primary and mixed-prime refinements remain
in the all-finite-quotient intersection of D.22.
The alternate certificate changes the initial homology
action in the same CM matrix pair, explaining exactly
why the rejection of the twenty-three selected actions
does not reject that matrix pair. The surviving alternate
action is recorded separately and is not substituted
silently for any original D.25 certificate.
\end{proof}
